% Matrix extension; no vector empirical superiority is inferred.
\section{A Matrix Form of Paid Calibration Value}
The scalar initialization argument relies on the fact that every program contracts all error directions by the same scalar.
For vector learners this property can fail, so replacing every program's initial error by its separate worst direction can produce an unattainable score.
The following reduction keeps one common unknown initialization across the whole observation space.
Its minimax interchange is a classical convex-game argument; the specific object is the finite-dimensional covariance problem for paid continuing programs.

\subsection{Information, affine continuation and integrability}
Let $e\in\mathbb R^d$ satisfy $e^\top P e\le r^2$, with public $P\succ0$ and $r>0$.
The unobserved $e$ is fixed across the experiment and does not affect the centered calibration-probe law.
Fix a public prior over a finite set of Gaussian noise models $j$.
For each model, the probe law $Y\sim\mathcal N(0,\Sigma_j)$ is nonsingular, the conditional common state $C\in\mathbb R^k$ has linear mean $m_j(y)$ and constant covariance, and the continuing Gaussian noise is jointly specified.
Both the paid probe schedule and every subsequent candidate program are frozen before observing $Y$.
Under model $j$ and program $a$, suppose the terminal error has conditional mean $A_{j,a}e+B_{j,a}m_j(y)$ and conditional covariance $V_{j,a}\succeq0$.
All matrices are known and deterministic, and their horizons include the probe charges, affine delay and withheld updates.
An implementable selector sees $Y$ and these public objects but not $e$ or $j$.
In particular the centered calibration observation must be obtainable without subtracting an unknown target.

Let $g$ be the mixture density and $p_j(y)$ its model posterior, and define
\begin{align*}
H_a(y)&=\sum_jp_j(y)A_{j,a}^\top A_{j,a},\\
Q_a(y)&=\sum_jp_j(y)\{\operatorname{tr}V_{j,a}+\|B_{j,a}m_j(y)\|^2\}.
\end{align*}
Both functions are even in $y$, $H_a$ is bounded, and $Q_a$ is integrable under $g$ because the catalogue is finite and Gaussian second moments are finite.
The posterior squared loss equals $e^\top H_a(y)e+Q_a(y)$ plus a term linear in $e$ and odd in $y$.
Consequently symmetrizing a selector cannot increase its maximum risk over the centrally symmetric error set.

\subsection{Exact covariance reduction}
Let $\Delta_{\rm even}$ be the measurable randomized even selectors over the finite catalogue, and put
\[
\mathcal S=\{S\succeq0:\operatorname{tr}(PS)\le r^2\}.
\]
Then the exact robust risk following the paid probe is
\begin{align}
&\inf_{\delta}\sup_{e^\top Pe\le r^2}R_e(\delta)\nonumber\\
&=\max_{S\in\mathcal S}\int\min_a\{\operatorname{tr}[H_a(y)S]+Q_a(y)\}g(y)\,dy.
\label{eq:matrix-covariance-envelope}
\end{align}
The minimax error-direction problem has therefore reduced to a concave maximization over a finite-dimensional initialization covariance.
The probe integration is still present; its numerical approximation requires its own accuracy justification.

To prove Eq.~\eqref{eq:matrix-covariance-envelope}, first symmetrize any selector, using $R_e(\delta^{\rm sym})=[R_e(\delta)+R_{-e}(\delta)]/2$.
For an even selector define $M_\delta=\mathbb E_g\sum_a\delta_a(Y)H_a(Y)$ and $c_\delta=\mathbb E_g\sum_a\delta_a(Y)Q_a(Y)$.
Its risk is $e^\top M_\delta e+c_\delta$.
The convex hull of $ee^\top$ over $e^\top Pe\le r^2$ is exactly $\mathcal S$: diagonalize $P^{1/2}SP^{1/2}$ and express it as a convex combination of rank-one radius-$r$ errors, with any unused mass at zero.
Hence the worst-case risk equals $\max_{S\in\mathcal S}\operatorname{tr}(M_\delta S)+c_\delta$.

The set $\mathcal S$ is compact and convex.
Regard each component of $\delta$ as an element of $L^\infty(g)$.
The probability-simplex constraints and evenness are weak-star closed, so $\Delta_{\rm even}$ is weak-star compact and convex.
The integrable functions $H_a,Q_a$ make the objective continuous and affine in each argument for the weak-star and finite-dimensional topologies.
The classical minimax interchange applies.
At fixed $S$, an even first minimizer of the finitely many even scores minimizes the integral, proving Eq.~\eqref{eq:matrix-covariance-envelope}.
Continuity and compactness give attainment in the primal and covariance-dual problems.

\subsection{A dual maximizer does not license arbitrary tie breaking}
There exists a saddle pair $(\delta^*,S^*)$; an arbitrary pointwise minimizer at a dual maximizer $S^*$ need not be a minimax selector.
Its tie randomization must also make $S^*$ a worst-case covariance for that selector.
For example take an irrelevant nonsingular centered Gaussian probe, zero continuation noise, $P=I_2$, $r=1$ and two programs with $A_0=\operatorname{diag}(1,0)$ and $A_1=\operatorname{diag}(0,1)$.
The dual objective is $\min(S_{11},S_{22})$, maximized at $S^*=I_2/2$ with value $1/2$.
Both actions tie everywhere at $S^*$.
Selecting the first action deterministically has worst-case risk one; mixing the two actions equally attains $1/2$.
Thus the scalar fixed tie rule cannot be transported to arbitrary matrix contractions without checking the saddle condition.
If the minimizing action is unique $g$-almost everywhere at a dual maximizer, the unique pointwise minimizer is the saddle selector because a primal optimizer exists and must minimize there.

\subsection{The strongest fixed comparison is itself a semidefinite program}
For the full-budget no-probe catalogue let $\overline H_a^{\rm full}$ be the prior average of $A_{j,a}^\top A_{j,a}$ and $\overline Q_a^{\rm full}$ its full-budget zero-initialization risk.
A fixed randomized program uses a probability vector $w\in\mathcal W=\{w\ge0:\sum_aw_a=1\}$ independent of the calibration observations and unknown initialization.
Its optimum worst-case risk $I_{\rm fix}$ is the value of
\begin{equation}
\begin{aligned}
\underset{w\in\mathcal W,\,t\in\mathbb R}{\operatorname{minimize}}\quad
&r^2t+\sum_aw_a\overline Q_a^{\rm full}\\
\text{subject to}\quad&\sum_aw_a\overline H_a^{\rm full}\preceq tP.
\end{aligned}
\label{eq:matrix-fixed-SDP}
\end{equation}
The constraint bounds the largest generalized initialization eigenvalue exactly, so this is a semidefinite program over a finite catalogue.
The example above proves why a comparison restricted to deterministic fixed programs can overstate a matrix selector's benefit.
Strict superiority to every fixed randomized programme in this catalogue is attainable after the specified paid probe if and only if the value in Eq.~\eqref{eq:matrix-covariance-envelope} is strictly less than $I_{\rm fix}$.
The reverse inequality is an impossibility certificate for this probe and catalogue, rather than a statement about every possible adaptive experiment.

With $P=I$ and scalar contractions $A_{j,a}=d_aI$, every unknown direction gives the same coefficient and the maximizing covariance can be taken on the radius boundary.
The formula reduces to the previously proved scalar even conditional-risk selector, and a fixed random mixture cannot improve its linear scalar endpoint objective.
The scalar simulations evaluate this reduction; the matrix characterization is analytical.
